\hypertarget{appendix-f744e41b95b4}{}
\section*{Unit directions and safe normalization}
\section{Unit directions and safe normalization}\label{appendix-f744e41b95b4-unit-directions-and-safe-normalization}

\textbf{Curriculum depth:} Introduced · \textbf{First seen:} Phase6 · \textbf{Latest development:} Phase6

\subsection{Definition and mechanism}\label{appendix-f744e41b95b4-definition-and-mechanism}

A unit direction is finite, nonzero input scaled to approximate length one. Scale by the largest absolute component before measuring to avoid unnecessary squared overflow or underflow. Zero has no direction; application noise thresholds are distinct from numerical existence.

For finite nonzero input, scaled squared length lies between \(1\) and \(3\). A vector of huge components can have a computable unit direction even when its original length does not fit.

\subsection{Prerequisites}\label{appendix-f744e41b95b4-prerequisites}

\begin{itemize}
\tightlist
\item
  \hyperref[appendix-49933872bb51]{Norms lengths and induced distances} --- Lengths and their domains are needed before unit directions or angular measurements.
\item
  \href{../01-Foundations-I/main.pdf\#nameddest=appendix-7f1e17dcafb1}{Robust norm and intermediate range (01-Foundations-I, appendix-7f1e17dcafb1)} --- Range-safe scaling avoids losing a valid direction to unnecessary squared intermediates.
\end{itemize}

\subsection{Used by}\label{appendix-f744e41b95b4-used-by}

\begin{itemize}
\item
  \hyperref[appendix-e569e60e18e6]{Robust vector angles} --- Builds on this mechanism in the connected geometry story.
\item
  \hyperref[appendix-e31893dec985]{Surface tangents normals and bitangents} --- Builds on this mechanism in the connected geometry story.
\item
  \hyperref[appendix-44e00863b4d5]{Reflection refraction and shading directions} --- Builds on this mechanism in the connected geometry story.
\item
  \hyperref[appendix-380a220e2529]{Coordinate-space types and validated invariants} --- Builds on this mechanism in the connected geometry story.
\item
  \hyperref[appendix-cac01ebeeb11]{Axial depth and Euclidean range}
\item
  \hyperref[appendix-fc99ad39f351]{Omnidirectional camera projections}
\item
  \hyperref[appendix-f1a6d43e80d1]{Screen-space ray generation and mouse picking} \#\# Recall and next depth
\end{itemize}

Reconstruct the bounds and distinguish rejection, a caller fallback, and a validated wrapper. Full production proofs, target-specific integration and learner mastery require further evidence.

\subsection{Source and navigation}\label{appendix-f744e41b95b4-source-and-navigation}

\href{Continuous\%20Notes/Computing\%20Foundations\%20-\%20From\%20Bits\%20to\%20C++\%20and\%20Systems\%20APIs\#A\%20direction\%20can\%20exist\%20even\%20when\%20the\%20naive\%20length\%20overflows}{Teaching narrative} · \href{../01-Foundations-I/main.pdf\#nameddest=appendix-fd090d982a15-phase-6---parts-501-600}{Coverage and source gap (01-Foundations-I, appendix-fd090d982a15-phase-6-\/-\/-parts-501-600)} · \href{../01-Foundations-I/main.pdf\#nameddest=appendix-a103e7846d21-phase-6---source-analysis-and-verification}{Evidence and limits (01-Foundations-I, appendix-a103e7846d21-phase-6-\/-\/-source-analysis-and-verification)} · \hyperref[appendix-4ad47bdd269a]{Dependency Map}
